\chapter{Blockchain: re-entrancy as a failed distributive law}
\label{ch:blockchain}

The DAO attack of 2016 drained roughly a third of a live fund through one structural
flaw: a contract called out to an external party \emph{before} updating its own
bookkeeping, and the callee re-entered the still-open call and spent against stale
state. This chapter argues that this failure mode is not a heuristic ``smell'' but a
computable invariant. A smart contract is a coalgebra for a polynomial functor;
running two contracts that share a resource is not composition along a fixed handoff
topology (which cannot fail) but the \emph{interleaving} of two, and the interleaving
of two directed containers is a Zappa--Sz\'ep product whose existence is governed by a
single degree-two cohomology class $[\omega]$. Unprotected re-entrancy is exactly
$[\omega]\neq 0$: the weld fails, and the failure is a wrong interleaving. On the
minimal faithful model the class equals a concrete one-bit cocycle, and that identity
is machine-checked in Lean~4.

\section{The idea: re-entrancy is a failed distributive law}

We state the thesis before building the apparatus, because the apparatus exists to make
this one sentence precise and checkable.

\begin{quotation}
\noindent\emph{Composing two shared-resource contracts requires a distributive law of
their handoff categories; that law is a Zappa--Sz\'ep weld $\mathcal{C}\bowtie\mathcal{D}$;
the weld exists iff a degree-two class $[\omega]\in\Hcoh(\mathrm{Sk}_{\mathcal C};\Delta)$
vanishes; and unprotected re-entrancy is precisely the case $[\omega]\neq 0$.}
\end{quotation}

The three moves --- contract as coalgebra, interaction as Zappa--Sz\'ep product,
obstruction as an $\Hcoh$ class --- are developed in the next three sections. The
payoff, in Section~\ref{sec:blk-machine}, is that on the smallest instance faithful to
the DAO pattern the class collapses to a single field element $\varepsilon\in\mathbb{F}_2$:
the bit recording whether the re-entrant branch \emph{mutates} the caller's shared
state. The identity $[\omega]=\varepsilon$ is then a finite fact settled by a proof
assistant with no room for hand-waving.

Two grades travel with this chapter, matching the two source strata. The core
obstruction computation --- that the interleaving obstruction is a degree-two class and
that on the minimal model it equals the token-mutation bit, machine-checked in Lean ---
is carried at \textsc{[peer-reviewed]} (referee: Rick). The Ethereum/EVM modelling that
instantiates it --- blocks as containers, contracts as coalgebras, the DAO as the $s_2$
branch --- is carried at \textsc{[computed]}: a translation exercise, not a theorem.

\begin{remark}[why the degree is two, not zero or one]
Multi-agent-systems work places cellular-sheaf obstructions of degree $H^0$ (consensus)
and $H^1$ (identifiability) on the agents' \emph{communication} graph, asking whether the
parties currently agree. The container obstruction sits one degree up, on the
\emph{handoff} category, and asks the prior question: can the parties be fused into one
consistent joint system at all? Re-entrancy is the smallest place where the answer is
no. This degree contrast is argued in the source obstruction note \cite{cl-h2edges}.
\end{remark}

\section{Smart contracts as polynomial coalgebras}
\label{sec:blk-coalg}

We recall the container model of a contract, at grade \textsc{[computed]}, following the
Ethereum-as-container development.

\begin{definition}[contract as a polynomial coalgebra]
A \emph{smart contract} with state set $\Sigma$, input type $I$ (function signatures
with arguments), and output type $O$ (return values and emitted events) is a coalgebra
for the polynomial functor
\[
  p(X)=(O\times X)^{I},
  \qquad
  \mathrm{step}\colon \Sigma \longrightarrow (O\times\Sigma)^{I}.
\]
Given a state $\sigma$ and an input $i$, the contract emits an output $o\in O$ and a new
state $\sigma'\in\Sigma$.
\end{definition}

As a container, $p(X)=(O\times X)^I$ has a single shape (the interface is fixed),
positions indexed by $I$, and at each position the value $O\times X$: an observable output
paired with a next state. This is a Mealy machine; had the output depended only on the
state we would have the Moore machine $p(X)=O\times X^I$.

\begin{example}[the Gateway marketplace]
\label{ex:blk-gateway}
The Gateway NFT marketplace contract is such a coalgebra. Its state $\Sigma$ is the tuple
of on-chain mappings (token data, listings, authorities, pending fees, the next-id
counter); its input $I$ is the disjoint union of signatures
$\mathsf{mint}+\mathsf{list}+\mathsf{buy}+\mathsf{delist}+\mathsf{reserveTokenId}+
\mathsf{withdrawFees}+\cdots$; its output $O$ ranges over success/revert, the events
$\mathsf{Minted},\mathsf{Listed},\mathsf{Sold},\mathsf{Transfer}$, and return values. On
input $\mathsf{buy}(\mathrm{tokenId}{=}42,\mathrm{Commercial})$ the step function verifies
the listing, distributes payment (fee to minter, remainder to seller), mints a license
token, emits events, and returns the updated state.
\end{example}

The blockchain itself supplies the ambient directed container into which contracts are
embedded: a block is the container $H\tri P$ separating header (shape) from transaction
slots (positions), and the chain of blocks is a directed container --- equivalently a
small category whose objects are blocks and whose morphisms are ancestor paths, the
free category on the parent graph \cite{AU16}. We record this only to locate contracts;
the obstruction lives one level down, in how two contracts \emph{interleave}.

\begin{remark}[the pivot from the source]
The Ethereum note observes that two contracts interacting is a composition of
polynomial coalgebras, with interface the composition product $p_A\tri p_B$ on $\Poly$,
and states --- at grade \textsc{[computed]} --- that ``re-entrancy is what happens when
the composition $p_A\tri p_B$ has a non-trivial distributive law with $p_B\tri p_A$,
[and] the re-entrancy guard is precisely the refusal to admit such a distributive law.''
The remainder of this chapter is the promotion of that sentence to a measured invariant:
the distributive law is a Zappa--Sz\'ep weld, and its failure is a degree-two class.
\end{remark}

\section{Composition, interleaving, and the Zappa--Sz\'ep weld}
\label{sec:blk-zs}

We now move to grade \textsc{[peer-reviewed]}. The vocabulary is Ahman--Uustalu's; the
identification of the interleaving obstruction with a cohomology class is the source
note's contribution.

\begin{definition}[directed container, handoff category]
A directed container equips a container $(S,P)$ with a root, a sub-position operation, and
a position-concatenation. By Ahman--Uustalu, directed containers are exactly small
categories: objects are shapes, morphisms out of $s$ are positions $P(s)$, identities are
roots, composition is concatenation \cite{AU16}. We call this category the \emph{handoff
category} and write $\tri$ for the composition product of directed containers.
\end{definition}

\begin{definition}[Zappa--Sz\'ep product]
Let $\mathcal K$ be a small category and $\mathcal D\subseteq\mathcal K$ a wide
subcategory (same objects). A \emph{Zappa--Sz\'ep product} $\mathcal K=\mathcal C\bowtie
\mathcal D$ is a second wide subcategory $\mathcal C$ such that every morphism of
$\mathcal K$ factors \emph{uniquely} as a $\mathcal C$-morphism followed by a
$\mathcal D$-morphism --- a strict factorization system, equivalently a distributive law
$\lambda\colon\mathcal D\mathcal C\Rightarrow\mathcal C\mathcal D$ satisfying the
matched-pair axioms. For directed containers this is the distributive-law composition of
Ahman--Uustalu \cite{AU16}; its data is a matched pair of mutual actions.
\end{definition}

The pivotal fact, absent from Ahman--Uustalu and supplied by the source, is that the
\emph{complementary factor $\mathcal C$ may or may not exist}. Composition along a single
fixed topology is a functor and cannot fail; the interleaving of two shared-resource
systems is a Zappa--Sz\'ep product, and its existence is exactly what the obstruction
measures. The identification is:
\[
  \mathcal C\bowtie\mathcal D \text{ exists}
  \qquad\Longleftrightarrow\qquad
  [\omega]=0\in \Hcoh(\mathrm{Sk}_{\mathcal C};\Delta),
\]
the degree-two cohomology of the handoff category's skeleton with coefficients in the
reply (direction) functor $\Delta$ \cite{cl-h2edges}. The cohomology itself is classical
(Baues--Wirsching); the delta over Ahman--Uustalu is the identification of the
interleaving composite with $\mathcal C\bowtie\mathcal D$ and hence of its obstruction
with this class. The class is a categorical holonomy: each partial handoff completes
freely, but the completions fail to glue into a globally serialized system by a
monodromy around the dispatch/return loop \cite{cl-emergent}.

% NEWBIB: H.-J. Baues, G. Wirsching, "Cohomology of small categories", J. Pure Appl. Algebra 38 (1985), 187--211.

\section{The obstruction is a degree-two class, and it is machine-checked}
\label{sec:blk-machine}

We now specialize to the minimal faithful model of the caller/callee pattern and compute
the class. Everything in this section is \textsc{[peer-reviewed]}, and its finite core is
machine-checked in Lean~4.

\subsection*{The parametrized supervisor--worker category}

Fix three roles: $\mathsf S$ (supervisor / caller), $\mathsf W$ (worker / callee
running), $\mathsf R$ (a return point). The supervisor carries a one-bit \emph{turn
token}; toggling it is an involution $\tau$ with $\tau^2=\id$, so
$\operatorname{End}(\mathsf S)=\{\id_{\mathsf S},\tau\}\cong\Z/2$. The generating handoffs are
\[
  p,\ p\tau\colon\mathsf S\to\mathsf W \quad(\text{dispatch in turn-state $0$ resp.\ $1$}),
  \qquad
  s,\ s_2\colon\mathsf W\to\mathsf R \quad(\text{two worker outcomes}),
\]
with composites $q,q\tau\colon\mathsf S\to\mathsf R$ and right $\Z/2$-action
$p\cdot\tau=p\tau$, $q\cdot\tau=q\tau$. The prescribed right factor is the
supervisor-internal (token) part
$\mathcal D=\{\id_{\mathsf S},\tau,\id_{\mathsf W},\id_{\mathsf R}\}$, a wide subcategory
whose only nontrivial vertex group is $G_{\mathsf S}\cong\Z/2$; the left factor
$\mathcal C$ --- the genuine dispatch/return part --- is what a Zappa--Sz\'ep product
would have to supply.

\begin{definition}[the family $\mathcal K_\varepsilon$]
\label{def:blk-family}
For $\varepsilon\in\Z/2$ let $\mathcal K_\varepsilon$ be the category above with the single
length-two composite fixed by
\[
  s\circ p = q,
  \qquad
  s_2\circ p = q\cdot\tau^{\varepsilon},
\]
where $s$ is ``the worker returns normally'' and $s_2$ is ``the worker re-enters'' ---
calls back into the supervisor before returning. Thus
$\varepsilon=[\,s_2\circ p = q\tau\,]$ records whether the re-entrant outcome
\emph{mutates} the token ($\varepsilon=1$, unprotected re-entry) or \emph{fixes} it
($\varepsilon=0$, a state-protected re-entry / a lock).
\end{definition}

Each $\mathcal K_\varepsilon$ is a genuine category with $\mathcal D$ wide, and the reading
is exact: re-entrancy --- the callee mutating shared caller state during a pending call ---
is precisely the $s_2$-branch dragging the token when $\varepsilon=1$ \cite{macbeth-reduction}.

\subsection*{The finite $\mathbb{F}_2$ complex}

The categorical obstruction of \cite{cl-h2edges} is, for $\mathcal K_\varepsilon$,
computed by an explicit orbit/cochain calculation in the analytic source
\cite{macbeth-reduction}. We import its output; the reduction \emph{to} the complex is a
cited pen-and-paper input, isolated here as three separable facts (each a future
formalisation target).

\begin{proposition}[the reduction; {\cite{macbeth-reduction}}]
\label{prop:blk-red}
Fix $\varepsilon\in\Z/2$. Then:
\begin{itemize}
\item[\textup{(i)}] $\mathcal K_\varepsilon$ satisfies the freeness condition
  \textnormal{(L)} (each hom-presheaf is a free right $\Delta$-set) and the vertex-group
  hypothesis ($\Delta$ a disjoint union of abelian vertex groups with trivial left action);
\item[\textup{(ii)}] the skeleton orbit category $\mathrm{Sk}_{\mathcal C}$ is the branch
  $\mathsf S\to\mathsf W\rightrightarrows\mathsf R$, with both middle arrows $[s],[s_2]$
  composing after $[p]$ to the single arrow $[q]$, \emph{independent of $\varepsilon$};
\item[\textup{(iii)}] the normalized (Baues--Wirsching) cochain complex is
  $C^1\cong\mathbb{F}_2^2$, $C^2\cong\mathbb{F}_2^2$, $C^3=0$, with coboundary
  $\delta^1 h=(h[p]-h[q],\,h[p]-h[q])$ and defect $2$-cocycle of the natural transversal
  $T=\{p,q,s,s_2\}$ equal to $\omega_T(\varepsilon)=(0,\varepsilon)$.
\end{itemize}
\end{proposition}

Everything after Proposition~\ref{prop:blk-red} is finite linear algebra over
$\mathbb{F}_2$, and it is this part that is formalised. Since $C^3=0$ we have $Z^2=C^2$.
The coboundary image is the diagonal,
\[
  B^2=\operatorname{im}\delta^1=\{(t,t):t\in\mathbb{F}_2\}=\langle(1,1)\rangle,
\]
so $\Hcoh=C^2/B^2\cong\mathbb{F}_2$, realised by the \emph{class map}
\[
  \varphi\colon\mathbb{F}_2^2\to\mathbb{F}_2,
  \qquad
  \varphi(a,b)=b-a=a\oplus b,
\]
whose kernel is exactly $B^2$. The defect cocycle is
$\omega_\varepsilon:=\omega_T(\varepsilon)=(0,\varepsilon)$, so
$\varphi(\omega_\varepsilon)=\varepsilon$: trivial for $\varepsilon=0$ and the nonzero
generator for $\varepsilon=1$. Chaining with the obstruction theorem gives the dichotomy
\[
  \mathcal C\bowtie\mathcal D \text{ exists}
  \iff [\omega_\varepsilon]=0
  \iff \varepsilon=0 .
\]

\begin{remark}[reading the class]
The two coordinates of $\omega$ are the two worker outcomes; the coboundary diagonal is the
freedom to relabel which token-state is ``canonical.'' The class is nonzero precisely
because the two outcomes disagree on that relabelling by one token-flip --- the re-entrant
branch's side effect. This is the categorical form of ``a flat bundle is trivial iff its
holonomy vanishes'': each partial handoff completes freely, but the completions fail to glue
into a global serialized agent by a single $\Z/2$ monodromy around the dispatch/return loop.
\end{remark}

\subsection*{What is machine-checked}

The finite $\mathbb{F}_2$ computation is transcribed into Lean~4 and machine-checked. The
development is sorry-free and uses no external library --- in particular no Mathlib ---
because the entire computation is finite, so the substantive facts close by \texttt{decide}
or case exhaustion. The model takes $\mathbb{F}_2=\texttt{Bool}$ with \texttt{xor} as
addition ($=$ subtraction in characteristic two), and defines $C^2$, $\delta^1$, $B^2$,
$\omega_T$, $\varphi$ as \texttt{C2}, \texttt{d1}, \texttt{InB2}, \texttt{omega},
\texttt{phi}.

\begin{theorem}[the certified statement; {\cite{macbeth-reduction}}]
\label{thm:blk-lean}
In \texttt{Containers/Reentrancy.lean}, sorry-free and without external libraries, the
following are machine-checked:
\begin{itemize}
\item \texttt{phi\_ker\_eq\_inB2} and \texttt{phi\_surjective}, together giving the
  isomorphism $\Hcoh=C^2/B^2\cong\mathbb{F}_2$ realised by $\varphi$;
\item \texttt{phi\_omega}: $\varphi(\omega_\varepsilon)=\varepsilon$ --- the class form of
  $[\omega]=\varepsilon$;
\item \texttt{omega\_inB2\_iff\_zero}: $\omega_\varepsilon\in B^2\iff\varepsilon=0$ --- the
  membership form, i.e.\ the joint agent $\mathcal C\bowtie\mathcal D$ exists iff the worker
  fixes the token.
\end{itemize}
Consequently, on the finite complex of Proposition~\ref{prop:blk-red}, the obstruction class
equals the token-mutation bit, and unprotected re-entrancy ($\varepsilon=1$) is the nonzero
generator of $\Hcoh\cong\Z/2$.
\end{theorem}

The proofs are as short as they read: \texttt{phi\_omega} is a two-case split closing by
\texttt{rfl}, \texttt{omega\_inB2\_iff\_zero} rewrites by $\ker\varphi=B^2$ then by
$\varphi(\omega_\varepsilon)=\varepsilon$, and \texttt{phi\_ker\_eq\_inB2} is a four-case
exhaustion over $(a,b)\in\mathbb{F}_2^2$. The brevity is the point: once the categorical
obstruction is reduced to this complex, $[\omega]=\varepsilon$ is a finite fact a proof
assistant settles with no room for hand-waving.

\subsection*{The verification boundary, drawn honestly}

The Lean development certifies the \emph{finite $\mathbb{F}_2$ class computation}: given the
complex of Proposition~\ref{prop:blk-red}, it proves $\Hcoh\cong\mathbb{F}_2$,
$\varphi(\omega_\varepsilon)=\varepsilon$, and the membership dichotomy. It does \emph{not}
certify the \emph{reduction} of the categorical obstruction to that complex --- that
$\mathcal K_\varepsilon$ satisfies (L) and the vertex-group hypothesis, that
$\mathrm{Sk}_{\mathcal C}$ is the stated branch, and that the transversal defect is
$(0,\varepsilon)$ (all of Proposition~\ref{prop:blk-red}) --- which is established by hand in
\cite{macbeth-reduction} and imported here as input. In one line: \emph{the last mile ---
from cochain complex to $[\omega]=\varepsilon$ --- is machine-checked; the first mile ---
from category to cochain complex --- is a cited pen-and-paper theorem.}

% TODO(lean): formalise the first mile --- the orbit category Sk_C, the freeness argument (L), and the transversal defect computation omega_T=(0,e) --- currently pen-and-paper in \cite{macbeth-reduction}.

\section{The EVM instantiation: the DAO and the token bit}
\label{sec:blk-evm}

We close by grounding the model in the EVM, at grade \textsc{[computed]}. The claim is that
$\mathcal K_\varepsilon$ is the DAO pattern stripped to its skeleton, and that the design
bit $\varepsilon$ is a familiar engineering quantity.

Read the three roles concretely. The supervisor $\mathsf S$ is the vulnerable contract
holding the shared resource --- in the DAO, the balance ledger. The dispatch $p$ is the
external call: \texttt{msg.sender.call\{value: \ldots\}("")}, made \emph{before} the
ledger is updated. The worker $\mathsf W$ is the callee, an attacker-controlled fallback
function. Its two outcomes are the crux:
\begin{itemize}
\item $s$ --- the worker returns normally, control comes back to $\mathsf S$, and the ledger
  is decremented as intended;
\item $s_2$ --- the worker \emph{re-enters}: its fallback calls back into $\mathsf S$'s
  withdraw entry point while the first call is still open and the ledger still reads the
  pre-withdrawal balance.
\end{itemize}
The token is the shared state whose consistency across the pending call is the whole
question. The bit $\varepsilon$ is whether the re-entrant branch $s_2$ mutates that state
relative to the normal branch $s$: for the DAO, $\varepsilon=1$, because the second
withdrawal succeeds against stale state and the composite $s_2\circ p$ lands in the
token-flipped composite $q\tau$ rather than the serialized $q$. By
Theorem~\ref{thm:blk-lean}, $[\omega]=\varepsilon=1\neq 0$: the two contract invocations
admit \emph{no} Zappa--Sz\'ep weld into a single well-defined serialized agent, and the
``composite contract'' the attacker exercises is a genuinely different --- wrongly
interleaved --- system.

\begin{remark}[the re-entrancy guard is a coboundary]
The standard mitigations are exactly the move $\varepsilon:1\mapsto 0$. The
checks-effects-interactions pattern updates the ledger \emph{before} the external call, so
the $s_2$ branch reads the already-decremented balance and $s_2\circ p=q$: the token is
fixed, $\varepsilon=0$. A mutex re-entrancy guard (\texttt{nonReentrant}) forbids the $s_2$
edge outright by refusing the second entry while the first is open --- in the source's
words, ``the re-entrancy guard is precisely the refusal to admit such a distributive law''
\cite{macbeth-reduction}. Either way the defect cocycle moves from $(0,1)$ to $(0,0)$, into
the coboundary $B^2=\langle(1,1)\rangle$, and the weld exists. The guard does not
\emph{detect} the obstruction; it \emph{annihilates} the class.
\end{remark}

\begin{remark}[scope and honesty]
Two limits are worth stating. First, $\mathcal K_\varepsilon$ is the minimal faithful model:
one supervisor, one worker, a one-bit token, a single length-two composite. Larger handoff
topologies --- several mutually re-entrant contracts, deeper call stacks --- carry a class in
a genuinely larger $\Hcoh$, and reducing them to a scalar is not claimed here. Second, the
EVM reading of this section is a \textsc{[computed]} translation: it asserts that the DAO's
call graph instantiates $\mathcal K_1$ and that guards realize $\varepsilon=0$, but it does
not formalise the EVM semantics. What is \textsc{[peer-reviewed]} and machine-checked is the
mathematical core: on $\mathcal K_\varepsilon$, unprotected re-entrancy is the nonzero
generator of $\Hcoh\cong\Z/2$, and the joint contract exists iff the callee fixes the shared
token.
\end{remark}

% NEWBIB: The DAO reentrancy exploit (June 2016) --- see e.g. P. Daian, "Analysis of the DAO exploit", Hacking Distributed, 2016; and the survey in N. Atzei, M. Bartoletti, T. Cimoli, "A survey of attacks on Ethereum smart contracts (SoK)", POST 2017, LNCS 10204, 164--186.

The through-line of the dossier's grant narrative appears here in its sharpest form:
compositional correctness has direct economic consequences, the correctness criterion is a
cohomology class, and on the canonical failure that class is not merely definable but
\emph{finite and machine-checked}. Re-entrancy is the smallest place where two systems
cannot be fused, and $[\omega]=\varepsilon$ is why.
